\section{Method}
\label{sec:method}

APT-Bench is intentionally broader than any single model. We nevertheless need a hierarchy-aware and novelty-aware reference architecture that can operate directly on APT profiles under the benchmark protocols. \ourmethod{} plays this role. It focuses on the genuine cell-level hierarchy, lineage $\rightarrow$ subtype, and leaves disease recognition to a separate lightweight front end.

\subsection{Cell Encoder}
\label{sec:method:backbone}

Each cell is represented by a standardized APT vector $x \in \mathbb{R}^{293}$. We embed feature identity and feature value separately. For aptamer index $i$, we form a token
\begin{equation}
t_i = \mathrm{Emb}(i) + \mathrm{MLP}_{\mathrm{val}}(x_i), \qquad i=1,\dots,293.
\end{equation}
A learnable \texttt{[CLS]} token is prepended to the resulting sequence and processed by a pre-norm Transformer encoder with 4 layers, 8 attention heads, hidden dimension 256, feed-forward dimension 512, GELU activations, and dropout 0.1. The final \texttt{[CLS]} representation
\begin{equation}
z = \mathrm{Encode}(x) \in \mathbb{R}^{256}
\end{equation}
serves as the shared cell representation for downstream tasks.

\subsection{Soft-Routed Hierarchical Cascade}
\label{sec:method:cascade}

Let $M=5$ be the number of coarse lineages and $K=\numsubtypes$ the number of fine subtypes. Each subtype $k$ belongs to exactly one parent lineage $m(k)$. A coarse head produces lineage logits
\begin{equation}
c = W_c z \in \mathbb{R}^{M},
\end{equation}
with routing distribution
\begin{equation}
p(m \mid x) = \mathrm{softmax}(c/\tau)_m,
\end{equation}
where $\tau=1$.

Instead of a hard cascade, which would force subtype prediction to follow only the top predicted lineage, \ourmethod{} uses a soft mixture of a global subtype path and lineage-specific experts:
\begin{equation}
s_k(x) =
\underbrace{g_k(z)}_{\text{global path}}
+
\alpha \, p(m(k)\mid x) \, e_{m(k),k}(z)
+
\gamma \log\!\big(\epsilon + (1-\epsilon)p(m(k)\mid x)\big),
\label{eq:cascade}
\end{equation}
Here $g(z)$ is a global subtype head, $e_{m,k}(z)$ is the expert associated with lineage $m$, $\alpha=1.0$, $\epsilon=0.1$, and $\gamma$ is ramped from 0 to 0.5 during early training. The global path prevents the model from collapsing to a hard hierarchy and allows recovery from coarse routing errors. Final predictions are
\begin{equation}
\hat{k} = \arg\max_k s_k(x),
\qquad
\hat{m} = \arg\max_m c_m.
\end{equation}

For analysis, we also compare soft routing with hard routing and oracle routing to separate errors caused by lineage uncertainty from those caused by subtype discrimination itself.

\subsection{Cross-Disease Contrastive Regularization}
\label{sec:method:contrastive}

We encourage subtype representations to align across disease contexts. Let
\begin{equation}
u = \frac{\mathrm{Proj}(z)}{\lVert \mathrm{Proj}(z) \rVert_2}
\end{equation}
be a normalized projection of the encoder representation. For anchor cell $i$, positives are cells with the same subtype but a different disease label. Same-subtype cells from the same disease are excluded because they do not test cross-disease invariance. Cells from different subtypes are treated as negatives, with additional emphasis on negatives that share the same coarse lineage. The contrastive loss weight is ramped from 0 to 0.2 after the early warmup period.

\subsection{Novelty Detection}
\label{sec:method:novelty}

Closed-set hierarchy prediction is insufficient for deployment because a new patient's cells may belong to a disease category absent from training. We therefore evaluate two novelty mechanisms.

\paragraph{Training-free scores.}
We first consider two standard training-free baselines. Given fine-level logits $s(x)$, we compute the maximum softmax probability novelty score
\begin{equation}
\mathrm{MSP}(x)
=
1-\max_k \mathrm{softmax}(s(x))_k
\end{equation}
\citep{hendrycks2016baseline}
and the energy score
\begin{equation}
\mathrm{Energy}(x)
=
-\log \sum_k \exp(s_k(x))
\end{equation}
\citep{liu2020energy}.
The same scores are computed from the lineage logits $c(x)$ to obtain coarse-level baselines.

\paragraph{Trained novelty head.}
For each LODO fold, we first train the encoder and cascade on the known diseases. We then freeze the representation and train a lightweight MLP novelty head on a proxy-novel disease selected from the known set:
\begin{equation}
\mathrm{NoveltyHead}: \mathbb{R}^{256} \rightarrow \mathbb{R}
\end{equation}
Cells from this proxy disease receive novelty label 1; cells from the remaining known diseases receive label 0. The true held-out disease is never used during either cascade training or novelty-head training. This nested design asks whether novelty learned from one simulated unseen disease transfers to a different truly unseen disease.

\subsection{Optimization and Model Scope}
\label{sec:method:training}

The classification cascade is trained jointly using
\begin{equation}
\mathcal{L}
=
\lambda_{\text{fine}}\mathcal{L}_{\text{fine}}
+
\lambda_{\text{coarse}}\mathcal{L}_{\text{coarse}}
+
w_{\text{contrast}}(t)\mathcal{L}_{\text{contrast}},
\end{equation}
where $\mathcal{L}_{\text{fine}}$ and $\mathcal{L}_{\text{coarse}}$ are cross-entropy losses with $\lambda_{\text{fine}}=1.0$ and $\lambda_{\text{coarse}}=0.5$. The contrastive weight increases to 0.2 according to the schedule in \S\ref{sec:method:contrastive}.

We optimize the cascade using AdamW with learning rate $10^{-4}$ and weight decay $10^{-5}$, gradient clipping at norm 1.0, batch size 512, and mixed-precision training. Models are trained for at most 50 epochs with early stopping based on validation fine-level macro-F1 with patience 10. The final configuration does not use class-balanced sampling, as our experiments found that it did not improve performance after correcting an implementation issue; additional details are provided in Appendix~\ref{sec:appendix:sampler}.

The novelty head is trained after the corresponding cascade checkpoint is fixed. We use AdamW with learning rate $10^{-3}$ and weight decay $10^{-4}$ for at most 30 epochs, with early-stopping patience 5. Disease recognition is not handled by a dedicated head inside \ourmethod{}. This is a deliberate modeling decision rather than a claim that disease is solved in a strong external sense: within-cohort disease prediction is already nearly saturated by simple classical baselines, while the architectural challenge addressed by \ourmethod{} is the unresolved cell-level hierarchy and novelty problem.
